% ============================================================================
\section{Introducción y marco metodológico}

Seguir de cerca la economía es fundamental para tomar decisiones públicas y empresariales, pero los responsables deben actuar casi siempre con cifras preliminares. En el caso del Producto Bruto Interno (PBI) mensual, la primera cifra que se publica no es definitiva: cambia en los meses siguientes a medida que llegan más encuestas de empresas, se actualizan registros tributarios y aduaneros, o se ajustan métodos de cálculo. La diferencia entre ese primer dato anunciado y su versión corregida meses después se conoce como revisión estadística.

Los estudios internacionales con datos en tiempo real muestran que estas correcciones no siempre son puro azar, sino que muchas veces siguen tendencias o patrones detectables con la información ya pública en el momento del anuncio \parencite{croushore2001,stark2010,croushore2011}. Aunque las técnicas modernas de factores y aprendizaje automático suelen usarse para estimar el crecimiento del mes actual antes de su publicación —lo que se llama \emph{nowcasting}— \parencite{giannone2008,banbura2013}, aquí se aborda un problema distinto: una vez que la primera cifra del PBI ya salió a la luz, ¿se puede predecir cuánto cambiará en el futuro?

Este trabajo respondió esa pregunta para el caso peruano y puso a prueba si un modelo de computación de reservorio cuántico (QRC) logra anticipar mejor esas revisiones que los métodos estadísticos y de aprendizaje automático tradicionales de complejidad equivalente. Desde el principio se evitó dar por sentada cualquier ventaja cuántica: el estudio se diseñó aceptando que los métodos clásicos simples podían igualar o incluso superar al enfoque cuántico.

El trabajo avanzó en cinco frentes ordenados. Primero, se revisó todo el archivo histórico de publicaciones del PBI peruano \parencite{cruz2025} para asegurar fechas confiables y separar los cambios de año base. Segundo, se analizaron los patrones estadísticos reales de las revisiones. Tercero, se pusieron a competir doce modelos distintos —estadísticos clásicos, aprendizaje automático, redes de reservorio clásicas y el reservorio cuántico— simulando pronósticos mes a mes sin conocer nunca el futuro. Cuarto, mediante doce escenarios simulados y cálculos de capacidad de detección, se verificó bajo qué condiciones teóricas podrían diferenciarse. Finalmente, se previó realizar una valoración económica práctica únicamente si el modelo cuántico demostraba una mejora clara y consistente, condición que no llegó a cumplirse.

El informe sigue este orden: la sección~2 explica el problema y la limpieza de los datos; la sección~3 repasa los estudios previos; la sección~4 presenta la justificación; y la sección~5 formula las preguntas centrales. La sección~6 explica los conceptos y modelos en términos claros; las secciones~7 y~8 plantean las hipótesis y objetivos; y la sección~9 detalla cómo se diseñaron las pruebas. La sección~10 expone los resultados encontrados; las secciones~11, 12 y~13 presentan la discusión, conclusiones y recomendaciones; y el documento cierra con las referencias bibliográficas (sección~14) y los anexos de respaldo~(A--L).
